\documentclass[11pt,a4paper]{article}

\usepackage[T1]{fontenc}
\usepackage[utf8]{inputenc}
\usepackage{lmodern}
\usepackage[a4paper,hmargin=26mm,vmargin=24mm]{geometry}
\usepackage{amsmath,amssymb,amsthm,mathtools}
\usepackage{stmaryrd}
\usepackage{microtype}
\usepackage{enumitem}
\PassOptionsToPackage{hyphens}{url}
\usepackage[hidelinks]{hyperref}
\usepackage[capitalise,noabbrev]{cleveref}
\crefname{equation}{equation}{equations}
\creflabelformat{equation}{(#2#1#3)}

\theoremstyle{plain}
\newtheorem{lemma}{Lemma}
\newtheorem*{proposition}{Proposition}
\theoremstyle{remark}
\newtheorem*{remark}{Remark}

\DeclareMathOperator{\dyn}{dyn}
\DeclareMathOperator{\ro}{ro}
\DeclareMathOperator{\Fb}{Fb}
\newcommand{\as}{\mathrm{as}}
\newcommand{\st}{\mathrm{st}}
\newcommand{\id}{\mathrm{id}}
\newcommand{\then}{\mathbin{\fatsemi}}
\newcommand{\MealyMod}{\mathbf{Mealy}/{\sim}}
\newcommand{\lbl}[1]{\nolinkurl{#1}}

\title{The feedback operator on $\MealyMod$\\ with separate dynamics and read-out}
\author{Working note for the ACT4E chapter on feedback}
\date{23 November 2026}

\begin{document}
\maketitle

\noindent
The example on Mealy machines in the section on feedback categories (Example 35.5 in the current build, equations (13)--(17)) is phrased with the dynamics and the read-out of a Mealy machine combined into one single function.
This note translates the example, and in particular its feedback operator, into the formulation with two separate functions, and proves that the translation is correct.
The translation is a short computation, but it brings out two things which the single-function formula hides: the feedback has a one-step delay, and the construction has a gap as soon as machines carry initial states.

\section{The example to be translated}

In the example, a morphism $A \to B$ of $\MealyMod$ is an equivalence class represented by a pair $\langle S, f \rangle$ of a set $S$ of states and a function
\begin{equation}
    f \colon A \times S \to B \times S .
\end{equation}
Given a morphism $[g] \colon A \times C \to B \times C$, represented by a function
\begin{equation}
    \label{eq:g}
    g \colon (A \times C) \times S \to (B \times C) \times S ,
\end{equation}
the example defines $\Fb_{A B}^{C}([g]) \colon A \to B$ to be the morphism represented by the function
\begin{equation}
    \label{eq:h}
    A \times (C \times S) \xrightarrow{\ \as^{-1}\ } (A \times C) \times S \xrightarrow{\ g\ } (B \times C) \times S \xrightarrow{\ \as\ } B \times (C \times S) .
\end{equation}
Here $\as$ denotes the re-bracketing function
\begin{equation}
    \as \colon (P \times Q) \times R \to P \times (Q \times R), \qquad \bigl((p, q), r\bigr) \mapsto \bigl(p, (q, r)\bigr),
\end{equation}
for the appropriate sets $P$, $Q$, $R$.
Functions are composed in diagrammatic order, written $\then$, as in the book.

\section{Setup and dictionary}

Following the book (\lbl{def:mealy-machine}), a Mealy machine from a set $U$ of inputs to a set $Y$ of outputs has a set $X$ of states, a dynamics and a read-out
\begin{equation}
    \dyn \colon U \times X \to X, \qquad \ro \colon U \times X \to Y,
\end{equation}
and an initial state.
The initial state is left aside until \cref{sec:initial}, because the example as written does not mention it.
The single-function form of the machine (\lbl{rem:mealy-single-function}) is
\begin{equation}
    \label{eq:single}
    F \colon U \times X \to Y \times X, \qquad F(u, x) = \bigl(\ro(u, x),\ \dyn(u, x)\bigr).
\end{equation}

\begin{lemma}[Dictionary]
    \label{lem:dictionary}
    Let $U$, $X$, $Y$ be sets.
    The assignment $(\dyn, \ro) \mapsto F$ of \cref{eq:single} is a bijection from the set of pairs of functions
    \begin{equation*}
        \bigl(\, U \times X \to X,\ \ U \times X \to Y \,\bigr)
    \end{equation*}
    to the set of functions $U \times X \to Y \times X$.
    Its inverse is $F \mapsto (F \then \pi_2,\ F \then \pi_1)$, where $\pi_1$ and $\pi_2$ are the projections of $Y \times X$.
\end{lemma}

\begin{proof}
    Let a pair $(\dyn, \ro)$ be given, and let $F$ be defined by \cref{eq:single}.
    Then $\pi_1(F(u, x)) = \ro(u, x)$ and $\pi_2(F(u, x)) = \dyn(u, x)$ for all $u \in U$ and $x \in X$, so $(F \then \pi_2, F \then \pi_1) = (\dyn, \ro)$.

    Conversely, let $F \colon U \times X \to Y \times X$ be given.
    Every element of $Y \times X$ is the pair of its two components, and hence
    \begin{equation*}
        F(u, x) = \bigl(\pi_1(F(u, x)),\ \pi_2(F(u, x))\bigr) \qquad \text{for all } u \in U,\ x \in X.
    \end{equation*}
    This says that $F$ is the single function associated by \cref{eq:single} to the pair $(F \then \pi_2, F \then \pi_1)$.
\end{proof}

So the phrase ``the two-function formulation of the machine given by $F$'' has exactly one meaning: its read-out is $F \then \pi_1$ and its dynamics is $F \then \pi_2$.

\section{The translated example}

Let $[g] \colon A \times C \to B \times C$ be represented by a machine with set of states $S$ and single function $g$ as in \cref{eq:g}.
By \cref{lem:dictionary}, applied with $U = A \times C$, $X = S$, and $Y = B \times C$, this is the same as the two functions
\begin{equation}
    \dyn_g \colon (A \times C) \times S \to S, \qquad \ro_g \colon (A \times C) \times S \to B \times C,
\end{equation}
with
\begin{equation}
    \label{eq:g-split}
    g(u, s) = \bigl(\ro_g(u, s),\ \dyn_g(u, s)\bigr) \qquad \text{for all } u \in A \times C,\ s \in S.
\end{equation}
Since the set of outputs is itself a product, the read-out splits further into its two components
\begin{equation}
    \ro_g^{B} \coloneqq \ro_g \then \pi_1 \colon (A \times C) \times S \to B,
    \qquad
    \ro_g^{C} \coloneqq \ro_g \then \pi_2 \colon (A \times C) \times S \to C,
\end{equation}
so that $\ro_g(u, s) = \bigl(\ro_g^{B}(u, s), \ro_g^{C}(u, s)\bigr)$.

\begin{proposition}
    The function of \cref{eq:h},
    \begin{equation*}
        h \coloneqq \as^{-1} \then g \then \as \colon A \times (C \times S) \to B \times (C \times S),
    \end{equation*}
    is the single function of the Mealy machine with set of inputs $A$, set of states $C \times S$, set of outputs $B$, read-out
    \begin{equation}
        \label{eq:ro-fb}
        \ro_{\Fb} \colon A \times (C \times S) \to B, \qquad
        \ro_{\Fb}\bigl(a, (c, s)\bigr) = \ro_g^{B}\bigl((a, c), s\bigr),
    \end{equation}
    and dynamics
    \begin{equation}
        \label{eq:dyn-fb}
        \dyn_{\Fb} \colon A \times (C \times S) \to C \times S, \qquad
        \dyn_{\Fb}\bigl(a, (c, s)\bigr) = \Bigl(\ro_g^{C}\bigl((a, c), s\bigr),\ \dyn_g\bigl((a, c), s\bigr)\Bigr).
    \end{equation}
    Moreover, $\ro_{\Fb}$ and $\dyn_{\Fb}$ are the only read-out and dynamics with this single function.
\end{proposition}

\begin{proof}
    Let $a \in A$, $c \in C$, and $s \in S$.
    Then
    \begin{align*}
        h\bigl(a, (c, s)\bigr)
         & = \as\Bigl(g\bigl(\as^{-1}(a, (c, s))\bigr)\Bigr) \\
         & = \as\Bigl(g\bigl((a, c), s\bigr)\Bigr) \\
         & = \as\Bigl(\ro_g\bigl((a, c), s\bigr),\ \dyn_g\bigl((a, c), s\bigr)\Bigr) \\
         & = \as\Bigl(\bigl(\ro_g^{B}((a, c), s),\ \ro_g^{C}((a, c), s)\bigr),\ \dyn_g\bigl((a, c), s\bigr)\Bigr) \\
         & = \Bigl(\ro_g^{B}\bigl((a, c), s\bigr),\ \bigl(\ro_g^{C}((a, c), s),\ \dyn_g((a, c), s)\bigr)\Bigr).
    \end{align*}
    The first equality is the definition of $h$, the second is the definition of $\as^{-1}$, the third is \cref{eq:g-split}, the fourth writes the element $\ro_g((a, c), s)$ of $B \times C$ as the pair of its components, and the fifth is the definition of $\as$.

    The result is an element of $B \times (C \times S)$.
    Its first component is the right-hand side of \cref{eq:ro-fb}, and its second component is the right-hand side of \cref{eq:dyn-fb}.
    Hence $h\bigl(a, (c, s)\bigr) = \bigl(\ro_{\Fb}(a, (c, s)),\ \dyn_{\Fb}(a, (c, s))\bigr)$, which is \cref{eq:single} for the machine in the statement.
    The uniqueness follows from \cref{lem:dictionary}, applied with $U = A$, $X = C \times S$, and $Y = B$.
\end{proof}

\paragraph{What this says.}
The $B$-component of the read-out is unchanged, except that the value $c$ is now read from the state instead of from the input.
The $C$-component of the old \emph{read-out} becomes part of the new \emph{dynamics}.
Concretely, if the machine is in the state $(c_{i-1}, s_{i-1})$ and receives the input $a_i$, then it returns the output $b_i$ and moves to the state $(c_i, s_i)$, where
\begin{equation}
    \begin{aligned}
        b_i & = \ro_g^{B}\bigl((a_i, c_{i-1}), s_{i-1}\bigr), \\
        c_i & = \ro_g^{C}\bigl((a_i, c_{i-1}), s_{i-1}\bigr), \\
        s_i & = \dyn_g\bigl((a_i, c_{i-1}), s_{i-1}\bigr).
    \end{aligned}
\end{equation}
So the value which $g$ emits on the $C$ wire at step $i$ is stored, and it is fed back to $g$ as its $C$ input at step $i + 1$.
This is feedback with a one-step delay.
The description ``shift $C$ into the state space'' of the single-function formulation does not show this; the two-function formulation does.

\section{Well-definedness on equivalence classes}

The feedback operator is defined by choosing a representative, so the translation is complete only once we have checked that it respects the equivalence relation.
I take ``congruence'' to mean isomorphism of Mealy machines, as in \lbl{def:MealyCat} and \lbl{def:mealy-machine-homomorphism}: a bijection between the sets of states which is compatible with the dynamics and with the read-outs.

\begin{lemma}
    \label{lem:well-defined}
    Let $g$ and $g'$ be Mealy machines from $A \times C$ to $B \times C$, with sets of states $S$ and $S'$, and let $\varphi \colon S \to S'$ be a bijection such that
    \begin{equation}
        \label{eq:iso}
        \varphi\bigl(\dyn_g(u, s)\bigr) = \dyn_{g'}\bigl(u, \varphi(s)\bigr)
        \qquad \text{and} \qquad
        \ro_{g'}\bigl(u, \varphi(s)\bigr) = \ro_g(u, s)
    \end{equation}
    for all $u \in A \times C$ and $s \in S$.
    Then the function
    \begin{equation*}
        \psi \coloneqq \id_C \times \varphi \colon C \times S \to C \times S', \qquad (c, s) \mapsto \bigl(c, \varphi(s)\bigr),
    \end{equation*}
    is an isomorphism of Mealy machines from the machine $(\dyn_{\Fb}, \ro_{\Fb})$ obtained from $g$ to the machine $(\dyn'_{\Fb}, \ro'_{\Fb})$ obtained from $g'$.
\end{lemma}

\begin{proof}
    The function $\psi$ is bijective because $\varphi$ is.
    Applying $\pi_1$ and $\pi_2$ to the second equation in \eqref{eq:iso} gives
    \begin{equation*}
        \ro_{g'}^{B}\bigl(u, \varphi(s)\bigr) = \ro_g^{B}(u, s)
        \qquad \text{and} \qquad
        \ro_{g'}^{C}\bigl(u, \varphi(s)\bigr) = \ro_g^{C}(u, s).
    \end{equation*}
    Let $a \in A$, $c \in C$, and $s \in S$.
    For the dynamics, we compute
    \begin{align*}
        \psi\bigl(\dyn_{\Fb}(a, (c, s))\bigr)
         & = \Bigl(\ro_g^{C}\bigl((a, c), s\bigr),\ \varphi\bigl(\dyn_g((a, c), s)\bigr)\Bigr) \\
         & = \Bigl(\ro_{g'}^{C}\bigl((a, c), \varphi(s)\bigr),\ \dyn_{g'}\bigl((a, c), \varphi(s)\bigr)\Bigr) \\
         & = \dyn'_{\Fb}\bigl(a, (c, \varphi(s))\bigr) \\
         & = \dyn'_{\Fb}\bigl(a, \psi(c, s)\bigr).
    \end{align*}
    For the read-outs, we compute
    \begin{equation*}
        \ro'_{\Fb}\bigl(a, \psi(c, s)\bigr)
        = \ro_{g'}^{B}\bigl((a, c), \varphi(s)\bigr)
        = \ro_g^{B}\bigl((a, c), s\bigr)
        = \ro_{\Fb}\bigl(a, (c, s)\bigr).
    \end{equation*}
    These are the two compatibility conditions for $\psi$.
\end{proof}

Hence $\Fb_{A B}^{C}([g])$ does not depend on the choice of the representative $g$.

\section{A gap: initial states}
\label{sec:initial}

The example describes a morphism of $\MealyMod$ as a pair $\langle S, f \rangle$.
In the book, however, a Mealy machine also carries an initial state $\st \in S$ (\lbl{def:mealy-machine}), and isomorphisms of Mealy machines are required to preserve it (\lbl{def:MealyCat}).
With initial states, the formula \eqref{eq:h} does not define a feedback operator.

The set of states of the new machine is $C \times S$, so its initial state must be a pair $(c_0, \st')$.
The choice $\st' = \st_g$ is natural, but nothing in the data determines the element $c_0 \in C$.
It is the initial value on the feedback wire, which the one-step delay described above requires.
Such an element can even fail to exist.
Let $C = \emptyset$.
Then $A \times C = \emptyset$ and $B \times C = \emptyset$, and for every set $S$ with a chosen element $\st$ there is a Mealy machine from $A \times C$ to $B \times C$ with set of states $S$, whose dynamics and read-out are empty functions.
But $C \times S = \emptyset$ has no element which could serve as the initial state.

There are three ways to repair this.

\begin{enumerate}[leftmargin=*]
    \item \emph{Drop the initial states} for this example: use machines $\langle S, \dyn, \ro \rangle$ up to bijections which are compatible with the dynamics and the read-out.
          Everything in this note then applies as it stands.
          As far as I know, this is the setting of the literature on feedback categories to which the to-do note above the example refers.
          It is, however, a different category from the one of \lbl{def:MealyCat}.
    \item \emph{Use pointed sets as objects}, so that every wire carries an initial value, and take $c_0$ to be the base point of $C$.
          The initial state of the new machine is then $(c_0, \st_g)$, and the function $\psi$ of \cref{lem:well-defined} preserves it whenever $\varphi$ preserves the initial states.
    \item \emph{Keep} \lbl{def:MealyCat} and accept that feedback needs $c_0$ as additional data.
          The operator then no longer has the type required by \lbl{def:feedback-cat}.
\end{enumerate}

\begin{remark}
    The definition of a feedback category (\lbl{def:feedback-cat}) asks for a \emph{strict} monoidal category, and the strict version of $\MealyMod$ in the book is the list-based one of \lbl{exa:MealyModL-strict-monoidal}.
    There the re-bracketing functions are identities, and the formulas \eqref{eq:ro-fb} and \eqref{eq:dyn-fb} simplify: on the set of concatenated tuples $\langle a, c, s \rangle$, the read-out of the new machine is $\ro_g^{B}$ and its dynamics is the pairing $\langle \ro_g^{C}, \dyn_g \rangle$, with no re-bracketing at all.
\end{remark}

\paragraph{Open decision.}
Which of the three options for the initial states to adopt.
This decides how the example is to be rewritten in the chapter.

\end{document}
